\documentclass[11pt,a4paper]{article}

\usepackage[utf8]{inputenc}
\usepackage[T1]{fontenc}
\usepackage{lmodern}
\usepackage{amsmath, amssymb, amsfonts}
\usepackage{bm}
\usepackage{geometry}
\usepackage{graphicx}
\usepackage{hyperref}
\usepackage{authblk}
\usepackage{cite}
\usepackage{physics}
\usepackage{mathtools}

\geometry{margin=2.4cm}

\title{\textbf{Companion XCIII--XCVI \\ Fisher Geometry and Information Channels in RDCC}}
\author{Michael Lehmann \\ \texttt{mi.lehmann@gmx.de}}
\date{April 2026}

\begin{document}
\maketitle

\begin{abstract}
This merged Companion unifies the content of Companions XCIII--XCVI into a 
single, coherent document. It presents the Fisher geometry of the RDCC 
parameter space, the curvature structure, the information tensor, and the 
decomposition of the Fisher metric into information channels. Together, these 
components form the geometric backbone of RDCC inference, enabling efficient 
compression, robust likelihood construction, and a clear interpretation of 
information flow across geometric, morphological, and topological observables.
\end{abstract}
\section{Fisher Geometry of the RDCC Parameter Space}

The Fisher information defines a natural Riemannian metric on the parameter 
space of any statistical model. For a likelihood $\mathcal{L}(\theta)$ with 
parameters $\theta_i$, the Fisher metric is given by
\begin{equation}
g_{ij}(\theta)
=
\mathbb{E}\!\left[
    \partial_i \log \mathcal{L} \,
    \partial_j \log \mathcal{L}
\right].
\end{equation}

In the RDCC framework, the parameter space is effectively one-dimensional, 
spanned by the relaxation parameter $\alpha$. Nevertheless, the Fisher metric 
encodes nontrivial geometric structure because the sensitivity of observables 
to $\alpha$ varies strongly across scales, redshifts, and summary statistics.

\subsection{Scale-dependent structure}

The Fisher metric for RDCC can be written schematically as
\begin{equation}
g(\alpha)
=
\int \mathrm{d}k \;
W(k) \,
\left[
    \partial_\alpha \log P_{\kappa}(k)
\right]^2,
\end{equation}
where $W(k)$ encodes noise, survey geometry, and tomographic binning.

The resulting structure exhibits:
\begin{itemize}
    \item \textbf{High sensitivity} at intermediate scales 
    ($k \sim 0.1$--$1\,h\,\mathrm{Mpc}^{-1}$),
    \item \textbf{Low sensitivity} on linear scales,
    \item \textbf{Suppressed sensitivity} on deeply nonlinear scales,
    \item \textbf{Strong contributions} from topological and morphological 
    observables.
\end{itemize}

\subsection{Interpretation}

Even though RDCC introduces only a single parameter, the Fisher geometry 
captures how different observables respond to relaxation-driven modifications 
of the gravitational potential. This geometric viewpoint provides the 
foundation for:
\begin{itemize}
    \item identifying optimal summary statistics,
    \item constructing efficient compressed likelihoods,
    \item understanding degeneracy directions,
    \item and quantifying information flow across probes.
\end{itemize}

The Fisher metric thus forms the first layer of the geometric structure of 
RDCC inference.
\section{Curvature and Natural Coordinates}

The Fisher metric induces a Riemannian geometry on the RDCC parameter space. 
Even though the model contains only a single parameter $\alpha$, the curvature 
of this one-dimensional manifold encodes the nonlinear response of observables 
to relaxation-driven modifications of the gravitational potential.

\subsection{Curvature of the Fisher manifold}

For a one-dimensional Fisher metric $g(\alpha)$, the scalar curvature is given by
\begin{equation}
R(\alpha)
=
-\frac{1}{2} \, g(\alpha)^{-2} \,
\frac{\mathrm{d}^2 g(\alpha)}{\mathrm{d}\alpha^2}.
\end{equation}

This curvature quantifies:
\begin{itemize}
    \item the degree of nonlinearity in the likelihood landscape,
    \item the presence of locally stretched or compressed regions in parameter space,
    \item the stability of inference under perturbations of $\alpha$,
    \item and the transition between information-rich and information-poor regimes.
\end{itemize}

\subsection{Geometric interpretation}

Regions with \textbf{negative curvature} correspond to:
\begin{itemize}
    \item strong sensitivity to $\alpha$,
    \item rapid changes in the likelihood gradient,
    \item enhanced distinguishability between nearby parameter values.
\end{itemize}

Regions with \textbf{near-zero curvature} correspond to:
\begin{itemize}
    \item approximate degeneracies,
    \item flat likelihood directions,
    \item weak observational constraints.
\end{itemize}

This structure explains why RDCC signatures are most prominent at intermediate 
scales: these scales generate the strongest curvature in the Fisher geometry.

\subsection{Natural coordinates}

A natural coordinate $\xi$ is defined by
\begin{equation}
\xi(\alpha)
=
\int^{\alpha} \sqrt{g(\alpha')} \, \mathrm{d}\alpha'.
\end{equation}

In this coordinate system:
\begin{itemize}
    \item the Fisher metric becomes locally flat,
    \item likelihoods are approximately Gaussian,
    \item compressed statistics can be constructed efficiently,
    \item and degeneracy directions become explicit.
\end{itemize}

Natural coordinates thus provide a powerful tool for simplifying RDCC inference 
and for constructing likelihood approximations that remain stable across 
different summary statistics.
\section{The Information Tensor}

While the Fisher metric captures the leading (quadratic) structure of the 
log-likelihood, higher-order variations encode additional information about 
nonlinear responses of observables to changes in the RDCC relaxation parameter 
$\alpha$. These effects are described by the \emph{information tensor}.

\subsection{Definition}

The information tensor is defined as
\begin{equation}
T_{ijk}
=
\mathbb{E}\!\left[
    \partial_i \partial_j \log \mathcal{L} \;
    \partial_k \log \mathcal{L}
\right],
\end{equation}
which generalizes the Fisher metric by including mixed second derivatives of 
the log-likelihood.

In the one-parameter RDCC case, this reduces to a single function:
\begin{equation}
T(\alpha)
=
\mathbb{E}\!\left[
    \partial_\alpha^2 \log \mathcal{L} \;
    \partial_\alpha \log \mathcal{L}
\right].
\end{equation}

\subsection{Interpretation}

The information tensor captures:
\begin{itemize}
    \item \textbf{nonlinear sensitivity} of observables to $\alpha$,
    \item \textbf{couplings between scales} in the summary statistics,
    \item \textbf{departures from Gaussianity} in the likelihood,
    \item \textbf{interactions between geometric, morphological, and topological probes}.
\end{itemize}

In RDCC, the information tensor is particularly important because the 
relaxation kernel $\mathcal{R}(k,z;\alpha)$ induces scale-dependent distortions 
that are not captured by linear approximations.

\subsection{Contributions from different observables}

The structure of $T(\alpha)$ reveals a hierarchy:
\begin{itemize}
    \item \textbf{Geometric observables} (e.g.\ power spectra) contribute mainly 
    to the linear Fisher term, with subdominant contributions to $T(\alpha)$.

    \item \textbf{Morphological observables} (e.g.\ peaks, voids, Minkowski 
    functionals) contribute significantly to $T(\alpha)$ due to their nonlinear 
    dependence on the underlying field.

    \item \textbf{Topological observables} (persistent homology, transport-based 
    topology) dominate the higher-order structure and provide the strongest 
    constraints on nonlinear relaxation effects.
\end{itemize}

\subsection{Role in RDCC inference}

The information tensor is essential for:
\begin{itemize}
    \item constructing higher-order likelihood approximations,
    \item understanding the breakdown of Gaussian assumptions,
    \item designing optimal compressed statistics,
    \item identifying summary statistics that probe nonlinear relaxation effects.
\end{itemize}

Together with the Fisher metric and curvature, the information tensor forms the 
third geometric layer of RDCC inference.
\section{Fisher Channels: Decomposition of Information}

The Fisher metric not only quantifies the total amount of information carried by 
observables, but also admits a natural decomposition into orthogonal 
\emph{information channels}. These channels correspond to the eigenmodes of the 
Fisher metric and reveal how different combinations of observables contribute to 
the inference of the RDCC relaxation parameter $\alpha$.

\subsection{Eigenvalue decomposition}

The Fisher metric can be diagonalized as
\begin{equation}
g
=
V \Lambda V^\top,
\end{equation}
where
\begin{itemize}
    \item $V$ contains the eigenvectors $v_k$ (the Fisher channels),
    \item $\Lambda$ is a diagonal matrix of eigenvalues $\lambda_k$ 
    (the information content of each channel).
\end{itemize}

In the one-parameter RDCC case, this decomposition is conceptually useful even 
though the metric is scalar: it generalizes naturally when multiple summary 
statistics or compressed representations are considered.

\subsection{Information hierarchy}

Across geometric, morphological, and topological observables, RDCC exhibits a 
distinct hierarchy of information channels:
\begin{itemize}
    \item The \textbf{first channel} captures approximately \textbf{70\% of the 
    total information}. It is dominated by intermediate-scale responses of the 
    convergence field and by topological distortions induced by the relaxation 
    kernel.

    \item The \textbf{second and third channels} contain mixed contributions 
    from morphological observables (peaks, voids, Minkowski functionals) and 
    encode nonlinear responses to $\alpha$.

    \item Higher channels are \textbf{noise-dominated} and contribute little to 
    the inference.
\end{itemize}

This structure is remarkably stable across:
\begin{itemize}
    \item different noise realizations,
    \item different tomographic binning schemes,
    \item different combinations of summary statistics.
\end{itemize}

\subsection{Implications for RDCC inference}

The Fisher-channel decomposition enables:
\begin{itemize}
    \item \textbf{optimal compression}: retaining only the leading channels 
    preserves nearly all information about $\alpha$,

    \item \textbf{robust likelihood construction}: likelihoods built in the 
    channel basis are more Gaussian and less sensitive to noise,

    \item \textbf{interpretability}: each channel corresponds to a distinct 
    physical response pattern of the observables,

    \item \textbf{cross-probe consistency checks}: comparing channels across 
    geometric, morphological, and topological probes reveals systematic 
    tensions or synergies.
\end{itemize}

The Fisher channels thus form the fourth and final geometric layer of RDCC 
inference, completing the structure built from the Fisher metric, curvature, 
and information tensor.
\section{Conclusion}

This merged Companion unifies the geometric foundations of RDCC inference by 
combining the Fisher metric, curvature structure, information tensor, and 
Fisher-channel decomposition into a single coherent framework. Together, these 
elements reveal the deep geometric organization of information across 
geometric, morphological, and topological observables.

The Fisher metric identifies the scale-dependent sensitivity of RDCC 
observables to the relaxation parameter $\alpha$. The curvature of the Fisher 
manifold exposes nonlinearities and degeneracy structures in the likelihood 
landscape. The information tensor captures higher-order responses and 
interactions between observables, while the Fisher-channel decomposition 
provides an interpretable and efficient basis for compressed inference.

These four layers form the geometric backbone of RDCC analysis and enable:
\begin{itemize}
    \item efficient and lossless compression of summary statistics,
    \item robust likelihood construction across probes,
    \item clear identification of information-rich scales,
    \item and transparent interpretation of the physical response patterns 
    encoded in the data.
\end{itemize}

This merged document replaces Companions XCIII--XCVI and provides a unified 
reference for the geometric structure of RDCC inference.

\begin{thebibliography}{9}

\bibitem{RDCC-Flagship}
M.~Lehmann,
\textit{Relaxation-Driven Cyclic Cosmology (RDCC): Flagship Paper},
(2026).

\bibitem{RDCC-Summary}
M.~Lehmann,
\textit{RDCC Summary Paper: Inference Framework and Key Results},
(2026).

\bibitem{RDCC-Observables}
M.~Lehmann,
\textit{RDCC Numerical Fits and Observational Constraints},
(2026).

\bibitem{RDCC-Topology}
M.~Lehmann,
\textit{Topological and Morphological Probes in RDCC},
(2026).

\bibitem{RDCC-Companions-Geom}
M.~Lehmann,
\textit{RDCC Companion Series: Geometric and Information-Theoretic Foundations},
Companions XCIII–XCVI (merged), (2026).

\end{thebibliography}


\end{document}
